% !TEX root = ../main.tex
\chapter{予備調査}
% 約5〜6ページ。データはそろっている

\section{先行発表(IPSJ)の実験と結果}
% 2台、タオルで囲った状態(CPU 面を向かい合わせ)、ファンなし
% 片側負荷で隣が約+12℃(data/legacy/expA, expB)、両側負荷では約84℃で頭打ち(expC)

\section{先行発表の課題}
% legacy/ipsj2026/linux_pc_sotuken/README.md と experiments/20261112_legacy_kf_check/README.md
% ・Δ秒先予測に未来の実測値(CPU使用率・隣接温度)を使っていた
% ・同定と評価が同じデータ。同じデータで設定を変えながら評価した
% ・CPU 使用率が瞬時値で記録され、負荷の係数 b が負(物理的に不自然)
% ・制御は実機に反映されていない(オフライン評価)。miss_count の定義の問題
% ・実験は約20分×1回、室温未記録、囲い(タオル)の再現性がない

\section{囲いなし・ラックでの熱干渉}
% experiments/20261113_openair_interference:囲いなしでは隣も離れたノードも +0.5℃ で差なし
% experiments/20261113_rack_interference:(結果待ち)

\section{本研究の方針}
% 干渉が出る条件と出ない条件を両方示す/干渉を受ける側も評価する/実際にファンを動かす/
% 室温を記録し、条件と手順を固定して複数回行う
